\section{Reading Comprehension Questions}


\begin{rem}
It is recommended that you attempt to complete {\em all} of the questions before checking your answers. 
As with the exercises throughout the book, it is also recommended that if you get one wrong, 
attempt to solve it again before reading the answer/solution in detail. You will learn a lot more that way!

Also, the solutions given are often just one possible answer (especially when answers involve coming up with an example). 
If your answer is different, 
you should be able to determine whether or not your answer is also correct. If you are not sure, please ask!
\end{rem}

\bigskip

\noindent
{\em From Section~\ref{section:logic}}

\begin{requ}
What is a proposition?
\begin{requanswer}
A proposition is a statement that is either true or false.
\end{requanswer}
\end{requ}

\begin{requ}
What are the six logical operators that were introduced in this chapter?
Draw a truth table for each.
\begin{requanswer}
The operators are {\em negation}, {\em or}, {\em and}, {\em exclusive-or}, {\em conditional}, and {\em biconditional}.
See Table~\ref{tab:not} and  Table~\ref{tab:2vars} for the truth tables.
\end{requanswer}
\end{requ}

\begin{requ}
Explain the difference between {\em (inclusive) or} and {\em exclusive or}.
\begin{requanswer}
$p\vee q$ is true if $p$ is true, $q$ is true, or both $p$ and $q$ are true, whereas
$p\oplus q$ is true if and only if exactly one of $p$ or $q$ is true. So the difference is
that is both $p$ and $q$ are true, $p\vee q$ is true, but $p\oplus a$ is false.
\end{requanswer}
\end{requ}

\begin{requ}
When is $p\rightarrow q$ true?
\begin{requanswer}
It is true unless $p$ is true and $q$ is false.
\end{requanswer}
\end{requ}

\begin{requ}
Draw a truth table for $(p\wedge q)\vee\neg p$.
\begin{requanswer}
Here is a truth table with some intermediate columns to help. As long as you have the same final column, yours is probably fine.

$\begin{array}{|cc||c|c|c|} \hline p & q & p  \wedge q & \neg p&(p \wedge q) \vee\neg p \\
\hline
T & T & T & F&T\\
T & F & F & F&F \\
F & T & F & T&T \\
F & F & F & T&T\\
\hline
\end{array}$
\end{requanswer}
\end{requ}

\begin{requ}
Draw a truth table for $(p\wedge q)\vee r$.
\begin{requanswer}
Here is a truth table with some intermediate columns to help. As long as you have the same final column, yours is probably fine.

$\begin{array}{|lll|l|l|}
\hline
p & q & r & p\wedge q& (p\wedge q)\vee r\\
\hline
T & T & T & T &T\\
T & T & F & T &T\\
T & F & T & F &T\\
T & F & F & F &F\\
F & T & T & F &T\\
F & T & F & F &F\\
F & F & T & F &T\\
F & F & F & F &F\\
\hline
\end{array}$
\end{requanswer}
\end{requ}

\begin{requ}
Can $p$ and $\neg p$ both be true? Explain.
\begin{requanswer}
By definition, no. If $p$ is true, $\neg p$ is false, and if $p$ is false, $\neg p$ is true.
\end{requanswer}
\end{requ}

\begin{requ}
If $p\vee q$ is true and $p$ is false, what can you say about $q$?
\begin{requanswer}
$q$ must be true since one of them has to be and $p$ is not.
\end{requanswer}
\end{requ}

\begin{requ}
If $p\vee q$ is true and $p$ is true, what can you say about $q$?
\begin{requanswer}
Nothing. It might be true, but it also might be false.
\end{requanswer}
\end{requ}

\begin{requ}
If $p\wedge q$ is true, what can you say about $p$ and/or $q$?
\begin{requanswer}
They are both true.
\end{requanswer}
\end{requ}

\begin{requ}
If $p\leftrightarrow q$ is true and $p$ is false, what can you say about $q$?
\begin{requanswer}
$q$ has to also be false since $p\leftrightarrow q$ implies they have the same truth value.
\end{requanswer}
\end{requ}

\begin{requ}
If $p\rightarrow q$ is true and $p$ is true, what can you say about $q$?
\begin{requanswer}
In this case $q$ also has to be true.
\end{requanswer}
\end{requ}

\begin{requ}
If $p\rightarrow q$ is true and $p$ is false, what can you say about $q$?
\begin{requanswer}
Since $p\rightarrow q$ is true whenever $p$ is false, we cannot say anything about $q$,
so it could be true or false.
\end{requanswer}
\end{requ}

\medskip

\noindent
{\em From Section~\ref{section:equiv}}

\begin{requ}
Can a proposition be a contingency and a tautology at the same time?
\begin{requanswer}
By definition, if a proposition is a contingency, then it is {\em not}
a tautology.
\end{requanswer}
\end{requ}

\begin{requ}
Is it possible for both a proposition and its negation to be true? Explain.
\begin{requanswer}
By definition, a proposition and its negation can {\em never} both be true.
If the proposition is true, its negation is false, and if the proposition is false, its negation is true.
\end{requanswer}
\end{requ}


\begin{requ}
Prove the domination laws. That is, prove that
\begin{enumerate}[(a)]
\item
$ p\vee T \ = T$ 
\item
 $p\wedge F = F$.
\end{enumerate}
\begin{requanswer}
(a) $p\vee T$ is true if and only if either $p$ or $T$ is true. Since clearly $T$ is true, $p\vee T$ is always true.
Therefore, $ p\vee T \ = T$.
Alternatively, you can give a truth table for $p\vee T$ and then comment something like 
``since the final column of the truth table is always true, $ p\vee T \ = T$.''
(b) We will do this one with a truth table: 
Notice that the column for $p\wedge F$ in the truth table below is always false. Therefore  $p\wedge F = F$.

$\begin{array}{|c||c|} \hline p & p\wedge F \\
\hline
T & F\\
F & F\\
\hline
\end{array}$
\end{requanswer}
\end{requ}			
			
			

\begin{requ}
Explain why $\neg p \wedge \neg q$ and $\neg(p\wedge q)$ are not logically equivalent.
\begin{requanswer}
You could draw a truth table and show that the columns for these two differ. However, there is any easier approach.
Notice that if $p$ is true and $q$ is false, $\neg p \wedge \neg q$ is false, but $\neg(p\wedge q)$ is true.
Therefore they are not equivalent.
\end{requanswer}
\end{requ}

\begin{requ}
Why is ``because that's not what DeMorgan's law says''
or ``because they look different'' not sufficient to prove that
$\neg p \wedge \neg q$ and $\neg(p\wedge q)$ are not logically equivalent.
\begin{requanswer}
We can't use the first reason because DeMorgan's law isn't the
only rule we have to determine whether or not two propositions
are equivalent. Perhaps they are equivalent by some other rule
(in this case they aren't, which you should know if you
answered the previous question).
The second reason is also not valid because 
every proposition is equivalent to many other 
propositions that look different than it, so just knowing that
they look different doesn't tell us anything. For instance,
DeMorgan's law says that $\neg p \wedge \neg q=\neg(p\vee q)$.
So even though those two don't look that same, they are equivalent.
\end{requanswer}
\end{requ}

\begin{requ}
What is an easy way to prove that two propositions are
not logically equivalent?
\begin{requanswer}
Find an assignment of true values to the variables
such that one is true and the other is false. That's all there is to it.
Although sometimes this is easier said than done!
\end{requanswer}
\end{requ}

\medskip

\noindent
{\em From Section~\ref{section:pred}}

\begin{requ}
What is a propositional function (or predicate)? Give an example.
\begin{requanswer}
A propositional function is a statement that contains one or more variables and depending on the values of the variable(s), 
the statement is true or false. In other words, a propositional function is a function whose outputs are propositions.
\end{requanswer}
\end{requ}

\begin{requ}
Does $\neg\forall x P(x)$ mean $P(x)$ is never true? If so, convince me. If not, what does it mean?
\begin{requanswer}
$\neg\forall x P(x)$ means that it is not the case that $P(x)$ is true for all values of $x$.
So it doesn't mean it is never true. It means that for one or more values of $x$ it is false.
That is, it means that it is not always true.
\end{requanswer}
\end{requ}

\begin{requ}
Does $\neg\exists x P(x)$ mean $P(x)$ is never true? If so, convince me. If not, what does it mean?
\begin{requanswer}
$\neg\exists x P(x)$ means that it is not the case that there is a value of $x$ for which $P(x)$ is true.
In other words, it is indeed saying that $P(x)$ is never true.
\end{requanswer}
\end{requ}

\begin{requ}
Give two equivalent (but different) ways of expressing $\forall x \neg \exists y Q(x,y)$.
\begin{requanswer}
The two obvious alternatives are 
$\neg\exists x  \exists y Q(x,y)$ and
$\forall x \forall y\neg  Q(x,y)$.
\end{requanswer}
\end{requ}

\begin{requ}
Give three equivalent (but different) ways of expressing $\neg \exists x  (x<0 \wedge x>0)$.
%(There are actually at least 4 reasonable equivalent expressions. Find them all if you can.)
%I'm not seeing the 4th. I'm not sure what I was thinking of. CAC, 5/21/21.
\begin{requanswer}
These are all equivalent: $ \forall x  \neg(x<0 \wedge x>0)$,
$ \forall x  (\neg(x<0) \vee  \neg(x>0))$, and
$ \forall x  (x>=0 \vee  x<=0)$.
\end{requanswer}
\end{requ}

\begin{requ}
Express the sentence ``Everybody hurts sometimes'' using predicates and quantifiers.
To get you started, let $H(x,y)=$``$x$ hurts at time $y$''.
\begin{requanswer}
$\forall x\exists y H(x,y)$.
\end{requanswer}
\end{requ}

% It's already an Problem.
%\begin{requ}
%Express the sentence ``No one ever is to blame'' using predicates and quantifiers.
%\end{requ}

\begin{requ}
Express the sentence ``Nothing ever changes, nothing ever stays the same'' using predicates and quantifiers.
Hint: You will need to define one or two predicates, depending on how you interpret the sentence and how clever you are.
\begin{requanswer}
Let $C(x,y)=$``$x$ changes at time $y$.'' 
Then $\neg C(x,y)=$``$x$ stays the same at time $y$.''
{\em If you did not define a predicate similar to this, stop reading now
and try to come up with the rest of the answer using $C(x,y)$ before continuing to read.}
Then the expression can be written as $\neg\exists x\exists y  C(x,y)\wedge \neg\exists x\exists y \neg C(x,y)$,
which of course can be simplified $\forall x\forall y \neg C(x,y) \wedge \forall x \forall y C(x,y)$
(but you would probably re-interpret this version in English as ``Everything stays the same, everything changes,''
which if you think about is it essentially saying the same thing.).
\end{requanswer}
\end{requ}

\begin{requ}
Let $P(x,y)=$``$x\leq y$'' and assume the universe of discourse is the set of integers.
\begin{enumerate}[(a)]
\item
Rephrase $\forall x\exists y P(x,y)$ in English.
\item
Rephrase $\exists x\forall y P(x,y)$ in English.
\item 
Do the statements in parts (a) and (b) seem to be saying the same thing? Explain.
\item
What is the truth value of $\forall x\exists y P(x,y)$?
\item
What is the truth value of $\exists x\forall y P(x,y)$?
\item
Hopefully you answered that one of the statements is true and the other is false
(If not, go back to the previous two questions and try again!).
Can you change the universe of discourse so that the two statements have the same truth values?
\item 
If you said no to the previous question, go back and try harder before continuing.
So, you can make them have the same truth value by changing the universe of discourse.  
Does that mean with this universe of discourse the statement are saying the same thing? 
(This is a subtle but important point, so if you are not totally confident in your answer, ask about this one!)
\end{enumerate}
\begin{requanswer}
(a) Given any integer, there is an integer that is at least as large as it.
(b) There is an integer such that every integer is at least as large as it is.
(c) They do not seem to be saying the same thing at all.
(d) True.
(e) False.
(f) If the universe of discourse is changed to positive integers, then both statements are true.
(The second one becomes true because every positive integer is at least as large as 1.)
(g) No. Just because two statements have the same truth value, that does not mean they are saying the same thing.
For instance, ``1+1=2" is true and ``$x^2\geq 0$'' is true, but they definitely are saying very different things. 
\end{requanswer}
\end{requ}

\medskip

\noindent
{\em From Section~\ref{section:normal}}

\begin{requ}
Let $p$, $q$ and $r$ be boolean variables. Which of the following are literals? 
$\neg p$, \, $\neg p \wedge r$, \, $q$, \, $\neg r$, \, $p\rightarrow r$, \, $r$.
\begin{requanswer}
$\neg p$, \, $q$, \, $\neg r$, \, $r$.
\end{requanswer}
\end{requ}

\begin{requ}
Let $p$, $q$ and $r$ be Boolean variables. Which of the following are conjunctive clauses? 
$\neg p$, \, $p\wedge r$, \, $q\vee \neg r$, \, $\neg p \wedge r$, \, $q$, \, $\neg r$, \, $p\rightarrow r$, \, $\neg(p\wedge q)$.
\begin{requanswer}
$\neg p$, \, $p\wedge r$, \, $\neg p \wedge r$, \, $q$,  \, $\neg r$.
\end{requanswer}
\end{requ}

\begin{requ}
Let $p$, $q$ and $r$ be Boolean variables. Which of the following are in disjunctive normal form? 
$\neg p$,\, 
$p\wedge r$, \, 
$q\vee \neg r$, \, 
$\neg p \wedge r$, \, 
$p\rightarrow r$, \, 
$\neg(p\wedge q)$,\, 
$(p\wedge q)\vee (q\wedge \neg r)\vee \neg p$,\, 
$(p\vee q)\wedge (q\vee \neg r)\wedge \neg p$,\, 
$(p\wedge r)\vee\neg(r\wedge \neg q)\vee (\neg p \wedge q)$,\, 
$(p\wedge \neg r\wedge q)\vee (\neg p\wedge r \wedge \neg q)\vee(p\wedge r \wedge q)\vee(\neg p\wedge \neg r \wedge \neg q)$
.
\begin{requanswer}
$\neg p$, \, 
$p\wedge r$, \, 
$q\vee \neg r$, \, 
$\neg p \wedge r$, \, 
$(p\wedge q)\vee (q\wedge \neg r)\vee \neg p$, \, 
$(p\wedge \neg r\wedge q)\vee (\neg p\wedge r \wedge \neg q)\vee(p\wedge r \wedge q)\vee(\neg p\wedge \neg r \wedge \neg q)$.
\end{requanswer}
\end{requ}

